\section*{Chapitre \ref{ch:b2:biomolecules} --- Biomolécules~: acides aminés, sucres, lipides et nucléotides}

\begin{solution}{exo:b2:biomolecules:1}
\begin{itemize}
  \item Entre pH 3 et 9, le \omterm{def:b2:biomolecules:amino-acid}{zwitterion} \ce{H3N+-CH(CH3)-COO-}.
  \item À pH 1, le cation \ce{H3N+-CH(CH3)-COOH}.
  \item À pH 12, l'anion \ce{H2N-CH(CH3)-COO-}.
\end{itemize}
\end{solution}

\begin{solution}{exo:b2:biomolecules:2}
Glycine~: $(2.35 + 9.78)/2 \approx 6.1$. Alanine~: $(2.35 + 9.87)/2 \approx 6.1$.
% ledger: pka:glycine1, pka:glycine2, pka:alanine1, pka:alanine2
\end{solution}

\begin{solution}{exo:b2:biomolecules:3}
\ce{H2N-CH(CH3)-CO-NH-CH2-COOH}~: l'alanine fournit l'extrémité N-terminale, la glycine l'extrémité
C-terminale. Gly--Ala, \ce{H2N-CH2-CO-NH-CH(CH3)-COOH}, en est un isomère de constitution~: le groupe
amino libre est sur la glycine.
\end{solution}

\begin{solution}{exo:b2:biomolecules:4}
\ce{NH2} à gauche~: L. Priorités \ce{NH2} > \ce{COOH} > \ce{CH2OH} > \ce{H}~: la L-sérine est
(\textit{S}).
\end{solution}

\begin{solution}{exo:b2:biomolecules:5}
Groupes~: ammonium N-terminal (vers 9), \omterm{def:b2:biomolecules:amino-acid}{chaîne latérale} de la lysine (10.7), \omterm{def:b2:biomolecules:amino-acid}{chaîne latérale} de
l'aspartate (3.9), groupe carboxylique C-terminal (vers 2). pH 1~: $+1 +1 + 0 + 0 = +2$. pH 7~: $+1 +1 -1 -1 = 0$. pH 12~:
$0 + 0 - 1 - 1 = -2$ (la \omterm{def:b2:biomolecules:amino-acid}{chaîne latérale} de la lysine, à 1.3 unité au-dessous de pH 12, est surtout neutre).
\end{solution}

\begin{solution}{exo:b2:biomolecules:6}
Le chlorure de thionyle et le chauffage attaqueraient aussi d'autres groupes et, surtout, le \omterm{def:b2:acyl-substitution:derivative}{chlorure
d'acyle} d'un acide aminé, très réactif, perd facilement sa configuration (l'hydrogène $\alpha$ devient
acide)~: le \omterm{def:b2:biomolecules:peptide}{peptide} serait en partie racémisé. Un \omterm{def:b2:biomolecules:peptide}{agent de couplage} active l'acide à température
ambiante, dans des conditions douces.
\end{solution}

\begin{solution}{exo:b2:biomolecules:7}
O5 lié à C2~: un cycle formé de C2, C3, C4, C5 et O, à cinq atomes (un furanose). En \omterm{def:b2:biomolecules:anomer}{projection de
Haworth}, O du cycle à l'arrière, C2 à droite avec \ce{CH2OH} (C1) vers le bas et \ce{OH} vers le haut
($\beta$)~; \ce{OH} de C3 vers le haut (à gauche
en Fischer), \ce{OH} de C4 vers le bas, C5 portant \ce{CH2OH} (C6) vers le haut.
\end{solution}

\begin{solution}{exo:b2:biomolecules:8}
$x_\alpha = (80.2 - 52.8)/(150.7 - 52.8) = 27.4/97.9 \approx 0.28$~: 28~\% d'$\alpha$, 72~\% de $\beta$.
\end{solution}

\begin{solution}{exo:b2:biomolecules:9}
Le maltose conserve un hémiacétal libre sur son second glucose, qui s'ouvre en aldéhyde~: il est
réducteur. Dans le saccharose, les deux \omterm{def:b2:biomolecules:anomer}{carbones anomères} sont engagés dans la \omterm{def:b2:biomolecules:glycoside}{liaison glycosidique}~:
pas d'hémiacétal, pas de pouvoir réducteur. Un acide dilué hydrolyse la \omterm{def:b2:biomolecules:glycoside}{liaison glycosidique}, un
acétal~: le glucose et le fructose sont libérés, et tous deux sont réducteurs.
\end{solution}

\begin{solution}{exo:b2:biomolecules:10}
La trioléine, \qty{885.45}{g/mol}, consomme trois \ce{KOH} (\qty{56.105}{g/mol})~: $3 \times 56.105/885.45 =
0.190$~g par g, soit un indice de \omterm{def:b2:acyl-substitution:saponification}{saponification} de 190. Des chaînes plus courtes signifient une masse
molaire plus petite pour les mêmes trois groupes ester~: plus de moles de corps gras par gramme, donc
plus de KOH.
% ledger: aw:C, aw:H, aw:O, aw:K
\end{solution}

\begin{solution}{exo:b2:biomolecules:11}
Protéger l'amine de la glycine sous forme de carbamate de \textit{tert}-butoxycarbonyle, et l'acide de
l'alanine sous forme d'ester méthylique. Coupler avec un carbodiimide. Retirer le carbamate par un acide
et l'ester par une base diluée (conditions orthogonales, si bien que chaque extrémité peut être libérée
seule)~: Gly--Ala.
\end{solution}

\begin{solution}{exo:b2:biomolecules:12}
5$'$-ACGCAT-3$'$ (lecture antiparallèle~: A s'apparie à T, G à C). Paires~: A--T, T--A, G--C, C--G, G--C, T--A~:
$3 \times 2 + 3 \times 3 = 15$ liaisons hydrogène.
\end{solution}

\begin{solution}{pb:b2:biomolecules:1}
\textbf{1.} L'acide L-aspartique, la L-phénylalanine et le méthanol.
\textbf{2.} Dans les deux cas, \ce{COOH} en haut, \omterm{def:b2:biomolecules:amino-acid}{chaîne latérale} en bas (\ce{CH2COOH}, \ce{CH2C6H5}), \ce{NH2}
à gauche.
\textbf{3.} (\textit{S}) pour les deux.
\textbf{4.} Deux stéréocentres~: quatre stéréoisomères.
\textbf{5.} \ce{H2N-CH(CH2COOH)-CO-NH-CH(CH2C6H5)-COOCH3}~: la \omterm{def:b2:biomolecules:peptide}{liaison peptidique} relie les deux
résidus~; l'ester est le \ce{COOCH3} de l'extrémité C-terminale.
\textbf{6.} Les récepteurs du goût sont chiraux~: ils fixent un stéréoisomère et non son image dans un miroir
ni ses diastéréoisomères, comme l'indiquait l'introduction.
\textbf{7.} Le \ce{NH3+} N-terminal et le \ce{COOH} de la \omterm{def:b2:biomolecules:amino-acid}{chaîne latérale} de l'acide aspartique~; l'acide
C-terminal est un ester méthylique, sans hydrogène acide.
\textbf{8.} À pH 2~: \ce{NH3+} ($+1$), \ce{COOH} de la \omterm{def:b2:biomolecules:amino-acid}{chaîne latérale} (0)~: $+1$.
\textbf{9.} À pH 7~: $+1 - 1 = 0$.
\textbf{10.} À pH 12~: $0 - 1 = -1$.
\textbf{11.} Entre les deux $\mathrm pK_a$ des groupes qui changent d'état~: $(3.9 + 9.9)/2 \approx 6.9$.
\textbf{12.} Dans le dipeptide, le carbone porteur de l'amine porte un amide au lieu d'un carboxylate~:
les groupes et les charges voisins changent, et les $\mathrm pK_a$ aussi (l'ammonium N-terminal d'un
\omterm{def:b2:biomolecules:peptide}{peptide} est nettement plus acide que celui d'un acide aminé libre). Le modèle donne l'allure, pas les
valeurs exactes.
\textbf{13.} L'acide et l'amine forment d'abord un sel~; le chauffage donnerait des mélanges (Asp--Asp,
chaînes plus longues, mauvais groupe carboxylique) et racémiserait les stéréocentres.
\textbf{14.} Son groupe amino et le groupe carboxylique de sa \omterm{def:b2:biomolecules:amino-acid}{chaîne latérale}.
\textbf{15.} Il transforme le \ce{OH} de l'acide en bon groupe partant, un dérivé activé que l'amine attaque
à température ambiante (substitution acyle).
\textbf{16.} Le $\beta$-dipeptide, lié par la \omterm{def:b2:biomolecules:amino-acid}{chaîne latérale}~: un isomère de l'aspartame, et non l'aspartame.
\textbf{17.} L'amine sous forme de carbamate de benzyloxycarbonyle et la \omterm{def:b2:biomolecules:amino-acid}{chaîne latérale} sous forme d'ester
benzylique~: tous deux sont retirés ensemble par hydrogénolyse sur palladium, qui laisse intact l'ester
méthylique (orthogonal à eux).
\textbf{18.} Un acide activé a un hydrogène $\alpha$ acide~; une base l'arrache et le stéréocentre se
racémise.
\textbf{19.} Protéger l'acide aspartique (carbamate sur N, ester benzylique sur la \omterm{def:b2:biomolecules:amino-acid}{chaîne latérale})~; activer
le $\alpha$-\ce{COOH} libre par l'\omterm{def:b2:biomolecules:peptide}{agent de couplage}~; ajouter l'ester méthylique de la L-phénylalanine~; hydrogénolyser.
\textbf{20.} L'ester méthylique et la \omterm{def:b2:biomolecules:peptide}{liaison peptidique}~; l'ester d'abord, un amide étant bien moins réactif.
\textbf{21.} Le dipeptide Asp--Phe (acide libre) et le méthanol.
\textbf{22.} La molécule sucrée est consommée~: ses produits d'hydrolyse ne la remplacent pas.
\textbf{23.} $14 \times 12.011 + 18 \times 1.008 + 2 \times 14.007 + 5 \times 15.999 =
\qty{294.307}{g/mol}$.
\textbf{24.} $0.180/294.307 = \qty{6.116e-4}{mol}$, soit \qty{0.612}{mmol}.
\textbf{25.} Un méthanol par aspartame~: $0.6116 \times 32.042 = \boldsymbol{\approx \qty{19.6}{mg}}$ de
méthanol.
\end{solution}
